%*******************************************************
% Abstract
%*******************************************************
\pdfbookmark[0]{Zusammenfassung}{Zusammenfassung}
\chapter*{Zusammenfassung}
In der vorliegenden Arbeit wird die Eignung lokaler Sprachmodelle für mobile Brain-Computer-Interface-Anwendungen (BCI) untersucht. Untersuchungsgegenstand sind drei lokal auf einem Android-Gerät ausgeführte Sprachmodelle, die in drei festgelegten Anwendungsszenarien anhand von Latenz, Ressourcenverbrauch, Stabilität und Ausgabequalität verglichen werden. In der bestehenden Applikation wurden Daten der Elektroenzephalographie (EEG) von einem tragbaren mobilen BCI-Gerät (Unicorn Hybrid Black) erfasst und anschließend visualisiert und/oder gespeichert. Solche Signale sind komplex und erfordern Fachwissen, um sie einwandfrei deuten zu können. Lokale Sprachmodelle könnten hier unterstützen, indem sie die aufbereiteten Messwerte in verständliche Sprache überführen. Das grundlegende Problem ist, dass solche Modelle oft auf ressourcenbeschränkten Geräten laufen und man daher die Tauglichkeit nicht garantieren kann. Daher ergibt sich die Forschungsfrage, inwieweit sich lokal auf Android-Geräten ausgeführte Sprachmodelle für die Verarbeitung von EEG-Daten im Rahmen mobiler Brain-Computer-Interface-Anwendungen eignen.
\vspace{\baselineskip}

\noindent Aus diesem Grund wurde die Anwendung um drei Kernkomponenten erweitert. Darunter fallen die Merkmalsextraktion, bei der die rohen EEG-Daten in eine für das Modell verständliche Form gebracht werden, die Anbindung lokaler Sprachmodelle (Gemma 3 1B, Gemma 2 2B, Qwen 2.5 1.5B) und die Implementierung einer Benutzeroberfläche. Zusätzlich wurde ein Benchmark implementiert, der für das automatisierte Messen der Latenz, des Ressourcenverbrauchs und der Stabilität der Modelle zuständig ist. Dabei ist die bestehende Anwendung selbst nicht Gegenstand der Bewertung, sondern dient als Messumgebung, in der die Modelle verglichen werden.
\vspace{\baselineskip}

\noindent Die Ergebnisse zeigen, dass die Verarbeitung der EEG-Daten durch alle drei getesteten Modelle stabil ist. Die Modellgröße bestimmt zwar die Zeit bis zum ersten Token (TTFT), Geschwindigkeit und Speicher, aber nicht automatisch Qualität oder Gesamtverbrauch. Keines der getesteten Modelle ist gleichzeitig in allen Bereichen überlegen. Die passende Modellwahl hängt schließlich davon ab, welcher Einsatzzweck verfolgt wird. Während Qwen 2.5 1.5B nicht überzeugend auftritt, sind die beiden Gemma-Modelle gut geeignet. Für schnelle Antworten auf ressourcenärmeren Geräten ist Gemma 3 1B die beste Wahl. Wenn Zeit keine Rolle spielt und die Antwort sowieso im Hintergrund generiert wird, ist Gemma 2 2B die bessere Wahl, da sie die beste Ausgabequalität liefert.
\vspace{\baselineskip}

\noindent Zusammenfassend lässt sich sagen, dass sich solche Modelle unter den untersuchten Bedingungen für die Verarbeitung von EEG-Daten eignen. Sie sind eine gute Ergänzung für die bestehende Anwendung und lassen viel Raum für weiterführende Untersuchungen.